\section{ANEXO A.~Ejemplos de vulnerabilidades en contratos
inteligentes}\label{anexo-a.-ejemplos-de-vulnerabilidades-en-contratos-inteligentes}

Este anexo presenta ejemplos simplificados de vulnerabilidades
representativas en contratos inteligentes desarrollados en Solidity. El
objetivo es ilustrar de forma práctica los principales tipos de
debilidades descritas en la sección 4.4, facilitando su comprensión y su
posterior detección mediante herramientas automáticas.

Cada ejemplo incluye: descripción, fragmento de código vulnerable,
impacto y estrategia de mitigación.

\subsection{ANEXO A.1. Vulnerabilidades técnicas de
ejecución}\label{anexo-a.1.-vulnerabilidades-tuxe9cnicas-de-ejecuciuxf3n}

\subsubsection{ANEXO A.1.1. Reentrancy}\label{anexo-a.1.1.-reentrancy}

La vulnerabilidad de
\href{https://www.cyfrin.io/blog/what-is-a-reentrancy-attack-solidity-smart-contracts}{\textbf{reentrancy}}
se produce cuando un contrato realiza una llamada externa antes de
actualizar su estado interno, permitiendo que el contrato receptor
reingrese en la función original en un estado inconsistente.

Código vulnerable

\begin{Shaded}
\begin{Highlighting}[]
\NormalTok{contract VulnerableBank \{}

    \FunctionTok{mapping}\NormalTok{(address }\KeywordTok{=\textgreater{}}\NormalTok{ uint256) }\KeywordTok{public}\NormalTok{ balances}\OperatorTok{;}

    \KeywordTok{function} \FunctionTok{deposit}\NormalTok{() }\KeywordTok{public}\NormalTok{ payable \{}

\NormalTok{        balances[msg}\OperatorTok{.}\AttributeTok{sender}\NormalTok{] }\OperatorTok{+=}\NormalTok{ msg}\OperatorTok{.}\AttributeTok{value}\OperatorTok{;}

\NormalTok{    \}}

    \KeywordTok{function} \FunctionTok{withdraw}\NormalTok{(uint256 amount) }\KeywordTok{public}\NormalTok{ \{}

        \PreprocessorTok{require}\NormalTok{(balances[msg}\OperatorTok{.}\AttributeTok{sender}\NormalTok{] }\OperatorTok{\textgreater{}=}\NormalTok{ amount)}\OperatorTok{;}

\FunctionTok{       }\NormalTok{ (bool success}\OperatorTok{,}\NormalTok{ ) }\OperatorTok{=}\NormalTok{ msg}\OperatorTok{.}\AttributeTok{sender}\OperatorTok{.}\AttributeTok{call}\NormalTok{\{}\DataTypeTok{value}\OperatorTok{:}\NormalTok{ amount\}(}\StringTok{""}\NormalTok{)}\OperatorTok{;}

        \PreprocessorTok{require}\NormalTok{(success)}\OperatorTok{;}

\NormalTok{        balances[msg}\OperatorTok{.}\AttributeTok{sender}\NormalTok{] }\OperatorTok{{-}=}\NormalTok{ amount}\OperatorTok{;}

\NormalTok{    \}}

\NormalTok{\}}
\end{Highlighting}
\end{Shaded}

\begin{itemize}
\tightlist
\item
  \textbf{Impacto}: Un atacante puede drenar fondos repitiendo la
  llamada antes de que el balance sea actualizado.
\item
  Mitigación:

  \begin{itemize}
  \tightlist
  \item
    Patrón
    \href{https://fravoll.github.io/solidity-patterns/checks_effects_interactions.html}{\emph{Checks-Effects-Interactions}}
  \item
    Uso de ReentrancyGuard
  \item
    Actualizar el estado antes de llamadas externas
  \end{itemize}
\end{itemize}

\subsubsection{ANEXO A.1.2. Integer Overflow /
Underflow}\label{anexo-a.1.2.-integer-overflow-underflow}

Errores aritméticos que provocan desbordamientos en operaciones enteras.
Aunque mitigados en Solidity \textgreater= 0.8, siguen siendo relevantes
en bloques unchecked.

Código vulnerable

\begin{Shaded}
\begin{Highlighting}[]
\KeywordTok{function} \FunctionTok{increment}\NormalTok{(uint256 x) }\KeywordTok{public}\NormalTok{ pure }\FunctionTok{returns}\NormalTok{ (uint256) \{}

\NormalTok{    unchecked \{}

        \ControlFlowTok{return}\NormalTok{ x }\OperatorTok{+} \DecValTok{1}\OperatorTok{;}

\NormalTok{    \}}

\NormalTok{\}}
\end{Highlighting}
\end{Shaded}

\begin{itemize}
\tightlist
\item
  \textbf{Impacto}: Puede alterar balances o condiciones lógicas
  críticas.
\item
  Mitigación:

  \begin{itemize}
  \tightlist
  \item
    Evitar unchecked salvo casos justificados
  \item
    Uso de validaciones explícitas
  \end{itemize}
\end{itemize}

\subsubsection{ANEXO A.1.3. Uso inseguro de
delegatecall}\label{anexo-a.1.3.-uso-inseguro-de-delegatecall}

La función delegatecall ejecuta código externo en el contexto de
almacenamiento del contrato llamador.

Código vulnerable

\begin{Shaded}
\begin{Highlighting}[]
\NormalTok{contract }\BuiltInTok{Proxy}\NormalTok{ \{}

\NormalTok{    address }\KeywordTok{public}\NormalTok{ implementation}\OperatorTok{;}

    \KeywordTok{function} \FunctionTok{execute}\NormalTok{(bytes memory data) }\KeywordTok{public}\NormalTok{ \{}

\FunctionTok{       }\NormalTok{ (bool success}\OperatorTok{,}\NormalTok{ ) }\OperatorTok{=}\NormalTok{ implementation}\OperatorTok{.}\FunctionTok{delegatecall}\NormalTok{(data)}\OperatorTok{;}

        \PreprocessorTok{require}\NormalTok{(success)}\OperatorTok{;}

\NormalTok{    \}}

\NormalTok{\}}
\end{Highlighting}
\end{Shaded}

\begin{itemize}
\tightlist
\item
  \textbf{Impacto}: Compromiso total del almacenamiento del contrato.
\item
  Mitigación:

  \begin{itemize}
  \tightlist
  \item
    Control estricto de la dirección implementation
  \item
    Uso de patrones proxy auditados
    (\href{https://eips.ethereum.org/EIPS/eip-1967}{EIP-1967},
    \href{https://docs.openzeppelin.com/contracts-stylus/uups-proxy}{UUPS})
  \end{itemize}
\end{itemize}

\subsubsection{ANEXO A.1.4. Denegación de servicio
(DoS)}\label{anexo-a.1.4.-denegaciuxf3n-de-servicio-dos}

Bloqueo de ejecución debido a fallos en llamadas externas o estructuras
no acotadas.

Código vulnerable

\begin{Shaded}
\begin{Highlighting}[]
\KeywordTok{function} \FunctionTok{payout}\NormalTok{(address[] memory recipients) }\KeywordTok{public}\NormalTok{ \{}

    \ControlFlowTok{for}\NormalTok{ (uint i }\OperatorTok{=} \DecValTok{0}\OperatorTok{;}\NormalTok{ i }\OperatorTok{\textless{}}\NormalTok{ recipients}\OperatorTok{.}\AttributeTok{length}\OperatorTok{;}\NormalTok{ i}\OperatorTok{++}\NormalTok{) \{}

        \FunctionTok{payable}\NormalTok{(recipients[i])}\OperatorTok{.}\FunctionTok{transfer}\NormalTok{(}\DecValTok{1}\NormalTok{ ether)}\OperatorTok{;}

\NormalTok{    \}}

\NormalTok{\}}
\end{Highlighting}
\end{Shaded}

\begin{itemize}
\tightlist
\item
  \textbf{Impacto}: Un solo fallo revierte toda la operación.
\item
  Mitigación

  \begin{itemize}
  \tightlist
  \item
    Uso de
    \href{https://medium.com/@markojauregui/the-pull-over-push-model-in-solidity-a-secure-pattern-for-fund-withdrawals-10c2e6628626}{patrón
    \emph{pull over push}}
  \item
    Evitar bucles dependientes de input externo
  \end{itemize}
\end{itemize}

\subsection{ANEXO A.2. Vulnerabilidades técnicas de control y
privilegios}\label{anexo-a.2.-vulnerabilidades-tuxe9cnicas-de-control-y-privilegios}

\subsubsection{ANEXO A.2.1. Falta de control de
acceso}\label{anexo-a.2.1.-falta-de-control-de-acceso}

Funciones críticas accesibles por cualquier usuario.

Código vulnerable

\begin{Shaded}
\begin{Highlighting}[]
\NormalTok{contract Ownable \{}

\NormalTok{    address }\KeywordTok{public}\NormalTok{ owner}\OperatorTok{;}

    \KeywordTok{function} \FunctionTok{withdrawAll}\NormalTok{() }\KeywordTok{public}\NormalTok{ \{}

        \FunctionTok{payable}\NormalTok{(msg}\OperatorTok{.}\AttributeTok{sender}\NormalTok{)}\OperatorTok{.}\FunctionTok{transfer}\NormalTok{(}\FunctionTok{address}\NormalTok{(}\KeywordTok{this}\NormalTok{)}\OperatorTok{.}\AttributeTok{balance}\NormalTok{)}\OperatorTok{;}

\NormalTok{    \}}

\NormalTok{\}}
\end{Highlighting}
\end{Shaded}

\begin{itemize}
\tightlist
\item
  \textbf{Impacto}: Pérdida total de fondos.
\item
  Mitigación:

  \begin{itemize}
  \tightlist
  \item
    Uso de onlyOwner
  \item
    Librerías como
    \href{https://docs.openzeppelin.com/contracts/5.x/access-control}{OpenZeppelin
    AccessControl}
  \end{itemize}
\end{itemize}

\subsubsection{ANEXO A.2.2. Uso de
tx.origin}\label{anexo-a.2.2.-uso-de-tx.origin}

Uso incorrecto de tx.origin para autenticación.

Código vulnerable

\begin{Shaded}
\begin{Highlighting}[]
\KeywordTok{function} \FunctionTok{withdraw}\NormalTok{() }\KeywordTok{public}\NormalTok{ \{}

    \PreprocessorTok{require}\NormalTok{(tx}\OperatorTok{.}\AttributeTok{origin} \OperatorTok{==}\NormalTok{ owner)}\OperatorTok{;}

    \FunctionTok{payable}\NormalTok{(msg}\OperatorTok{.}\AttributeTok{sender}\NormalTok{)}\OperatorTok{.}\FunctionTok{transfer}\NormalTok{(}\FunctionTok{address}\NormalTok{(}\KeywordTok{this}\NormalTok{)}\OperatorTok{.}\AttributeTok{balance}\NormalTok{)}\OperatorTok{;}

\NormalTok{\}}
\end{Highlighting}
\end{Shaded}

\begin{itemize}
\tightlist
\item
  \textbf{Impacto}: Ataques mediante contratos intermediarios.
\item
  \textbf{Mitigación}: Usar msg.sender para autenticación
\end{itemize}

\subsubsection{ANEXO A.2.3. Inicialización insegura (contratos
upgradeables)}\label{anexo-a.2.3.-inicializaciuxf3n-insegura-contratos-upgradeables}

En contratos upgradeables, la inicialización se realiza mediante
funciones externas (initialize) en lugar de constructores. Si no están
protegidas, cualquier usuario puede ejecutarlas y asumir el control del
contrato.

Código vulnerable

\begin{Shaded}
\begin{Highlighting}[]
\KeywordTok{function} \FunctionTok{initialize}\NormalTok{(address \_owner) }\KeywordTok{public}\NormalTok{ \{}

\NormalTok{    owner }\OperatorTok{=}\NormalTok{ \_owner}\OperatorTok{;}

\NormalTok{\}}
\end{Highlighting}
\end{Shaded}

\begin{itemize}
\tightlist
\item
  \textbf{Impacto}: Un atacante puede inicializar el contrato antes que
  el legítimo propietario.
\item
  Mitigación:

  \begin{itemize}
  \tightlist
  \item
    Uso de
    \href{https://docs.openzeppelin.com/upgrades-plugins/writing-upgradeable}{initializer
    (OpenZeppelin)}
  \item
    Bloqueo de inicialización tras ejecución
  \end{itemize}
\end{itemize}

\subsection{ANEXO A.3. Vulnerabilidades económicas y del
entorno}\label{anexo-a.3.-vulnerabilidades-econuxf3micas-y-del-entorno}

\subsubsection{ANEXO A.3.1. Front-running /
MEV}\label{anexo-a.3.1.-front-running-mev}

Un atacante observa la mempool y ejecuta transacciones antes que la
víctimANEXO A.

Código vulnerable

\begin{Shaded}
\begin{Highlighting}[]
\KeywordTok{function} \FunctionTok{buy}\NormalTok{(uint price) }\KeywordTok{public}\NormalTok{ \{}

    \PreprocessorTok{require}\NormalTok{(price }\OperatorTok{==}\NormalTok{ currentPrice)}\OperatorTok{;}

    \CommentTok{// compra}

\NormalTok{\}}
\end{Highlighting}
\end{Shaded}

\begin{itemize}
\tightlist
\item
  \textbf{Impacto}: Manipulación de operaciones (arbitraje,
  liquidaciones, subastas).
\item
  Mitigación:

  \begin{itemize}
  \tightlist
  \item
    \href{https://medium.com/coinmonks/commit-reveal-scheme-in-solidity-c06eba4091bb}{\emph{Commit-reveal}}
  \item
    Subastas ciegas
  \item
    Uso de \emph{relayers} privados
  \end{itemize}
\end{itemize}

\subsubsection{ANEXO A.3.2. Dependencia de
oráculos}\label{anexo-a.3.2.-dependencia-de-oruxe1culos}

Uso de datos externos manipulables.

Código vulnerable

\begin{Shaded}
\begin{Highlighting}[]
\KeywordTok{function} \FunctionTok{getPrice}\NormalTok{() }\KeywordTok{public}\NormalTok{ view }\FunctionTok{returns}\NormalTok{ (uint) \{}

    \ControlFlowTok{return}\NormalTok{ externalOracle}\OperatorTok{.}\FunctionTok{price}\NormalTok{()}\OperatorTok{;}

\NormalTok{\}}
\end{Highlighting}
\end{Shaded}

\begin{itemize}
\tightlist
\item
  \textbf{Impacto}: Manipulación de precios en DeFi.
\item
  Mitigación:

  \begin{itemize}
  \tightlist
  \item
    Oráculos descentralizados (ej.
    \href{https://chain.link/}{Chainlink})
  \item
    Promedios temporales
    (\href{https://www.binance.com/es-MX/support/faq/detail/80655cc54d8a4b2bb8ea097001844fd1}{TWAP})
  \end{itemize}
\end{itemize}

\subsubsection{ANEXO A.3.3. Uso de
block.timestamp}\label{anexo-a.3.3.-uso-de-block.timestamp}

El uso de block.timestamp como fuente de aleatoriedad o para decisiones
críticas es inseguro, ya que su valor puede ser parcialmente manipulado
por mineros o validadores dentro de ciertos límites.

Código vulnerable

\begin{Shaded}
\begin{Highlighting}[]
\KeywordTok{function} \FunctionTok{random}\NormalTok{() }\KeywordTok{public}\NormalTok{ view }\FunctionTok{returns}\NormalTok{ (uint) \{}

    \ControlFlowTok{return} \FunctionTok{uint}\NormalTok{(}\FunctionTok{keccak256}\NormalTok{(abi}\OperatorTok{.}\FunctionTok{encodePacked}\NormalTok{(block}\OperatorTok{.}\AttributeTok{timestamp}\NormalTok{)))}\OperatorTok{;}

\NormalTok{\}}
\end{Highlighting}
\end{Shaded}

\begin{itemize}
\tightlist
\item
  \textbf{Impacto}: Resultados predecibles o manipulables por
  mineros/validadores.
\item
  Mitigación:

  \begin{itemize}
  \tightlist
  \item
    VRF
    (\href{https://chain.link/education-hub/verifiable-random-function-vrf}{Verifiable
    Random Functions})
  \item
    Fuentes externas verificables
  \end{itemize}
\end{itemize}

\subsection{ANEXO A.4. Errores lógicos de
negocio}\label{anexo-a.4.-errores-luxf3gicos-de-negocio}

\subsubsection{ANEXO A.4.1. Error en cálculo de
balances}\label{anexo-a.4.1.-error-en-cuxe1lculo-de-balances}

Errores en operadores lógicos o condiciones de validación pueden
provocar inconsistencias en la gestión de balances, especialmente en
casos límite donde las condiciones no cubren todos los escenarios
posibles.

Código vulnerable

\begin{Shaded}
\begin{Highlighting}[]
\KeywordTok{function} \FunctionTok{withdraw}\NormalTok{(uint amount) }\KeywordTok{public}\NormalTok{ \{}

    \PreprocessorTok{require}\NormalTok{(balances[msg}\OperatorTok{.}\AttributeTok{sender}\NormalTok{] }\OperatorTok{\textgreater{}}\NormalTok{ amount)}\OperatorTok{;}

\NormalTok{    balances[msg}\OperatorTok{.}\AttributeTok{sender}\NormalTok{] }\OperatorTok{{-}=}\NormalTok{ amount}\OperatorTok{;}

\NormalTok{\}}
\end{Highlighting}
\end{Shaded}

\begin{itemize}
\tightlist
\item
  \textbf{Impacto}: Comportamiento incorrecto en condiciones límite.
\item
  Mitigación:

  \begin{itemize}
  \tightlist
  \item
    Uso de \textgreater{} en lugar de \textgreater=.
  \end{itemize}
\end{itemize}

\subsubsection{ANEXO A.4.2. Distribución incorrecta de
recompensas}\label{anexo-a.4.2.-distribuciuxf3n-incorrecta-de-recompensas}

La lógica de distribución puede introducir errores debido a divisiones
enteras o falta de gestión de restos, provocando pérdida de precisión y
fondos no asignados correctamente.

Código vulnerable

\begin{Shaded}
\begin{Highlighting}[]
\KeywordTok{function} \FunctionTok{distribute}\NormalTok{() }\KeywordTok{public}\NormalTok{ \{}

\NormalTok{    uint reward }\OperatorTok{=}\NormalTok{ total }\OperatorTok{/}\NormalTok{ users}\OperatorTok{.}\AttributeTok{length}\OperatorTok{;}

    \ControlFlowTok{for}\NormalTok{ (uint i }\OperatorTok{=} \DecValTok{0}\OperatorTok{;}\NormalTok{ i }\OperatorTok{\textless{}}\NormalTok{ users}\OperatorTok{.}\AttributeTok{length}\OperatorTok{;}\NormalTok{ i}\OperatorTok{++}\NormalTok{) \{}

\NormalTok{        balances[users[i]] }\OperatorTok{+=}\NormalTok{ reward}\OperatorTok{;}

\NormalTok{    \}}

\NormalTok{\}}
\end{Highlighting}
\end{Shaded}

\begin{itemize}
\tightlist
\item
  Impacto:

  \begin{itemize}
  \tightlist
  \item
    Pérdida de fondos debido a errores de redondeo (truncamiento en
    división entera)
  \item
    Acumulación de saldo no distribuido en el contrato
  \item
    Distribuciones injustas entre usuarios
  \item
    Posibles vectores de explotación si un atacante manipula el número
    de participantes
  \end{itemize}
\item
  Mitigación:

  \begin{itemize}
  \tightlist
  \item
    Uso de patrones de distribución que gestionen residuos (por ejemplo,
    acumuladores o ``\emph{remainder handling}'')
  \item
    Empleo de mayor precisión mediante escalado
    (\href{https://rareskills.io/post/solidity-fixed-point}{\emph{fixed-point
    arithmetic}})
  \item
    Validación de invariantes económicas (la suma distribuida debe
    coincidir con el total)
  \item
    Testing específico de casos límite (número de usuarios, valores
    pequeños, etc.)
  \end{itemize}
\end{itemize}

\subsubsection{ANEXO A.4.3. Estados
inconsistentes}\label{anexo-a.4.3.-estados-inconsistentes}

La falta de control adecuado sobre las transiciones de estado puede
permitir la ejecución de funciones en condiciones no válidas, generando
comportamientos inconsistentes en el contrato.

Código vulnerable

\begin{Shaded}
\begin{Highlighting}[]
\KeywordTok{enum}\NormalTok{ State \{ Open}\OperatorTok{,}\NormalTok{ Closed \}}

\NormalTok{State }\KeywordTok{public}\NormalTok{ state}\OperatorTok{;}

\KeywordTok{function} \FunctionTok{close}\NormalTok{() }\KeywordTok{public}\NormalTok{ \{}

\NormalTok{    state }\OperatorTok{=}\NormalTok{ State}\OperatorTok{.}\AttributeTok{Closed}\OperatorTok{;}

\NormalTok{\}}

\KeywordTok{function} \FunctionTok{bid}\NormalTok{() }\KeywordTok{public}\NormalTok{ payable \{}

    \PreprocessorTok{require}\NormalTok{(state }\OperatorTok{==}\NormalTok{ State}\OperatorTok{.}\AttributeTok{Open}\NormalTok{)}\OperatorTok{;}

\NormalTok{\}}
\end{Highlighting}
\end{Shaded}

\begin{itemize}
\tightlist
\item
  Impacto:

  \begin{itemize}
  \tightlist
  \item
    Ejecución de funciones en estados no válidos
  \item
    Comportamiento inesperado del contrato
  \item
    Bloqueo o bypass de lógica de negocio
  \item
    Posible explotación combinada con otras vulnerabilidades (por
    ejemplo, \emph{front-running} o \emph{reentrancy}
  \end{itemize}
\item
  Mitigación:

  \begin{itemize}
  \tightlist
  \item
    Implementación de máquinas de estados explícitas y completas
  \item
    Uso de modificadores para validar estado (inState(State.Open))
  \item
    Restricción de transiciones de estado válidas
  \item
    Aplicación de patrones
    \href{https://fravoll.github.io/solidity-patterns/state_machine.html}{\emph{state
    machines}}
  \end{itemize}
\end{itemize}

\subsection{ANEXO A.5. Resumen de
vulnerabilidades}\label{anexo-a.5.-resumen-de-vulnerabilidades}

{\def\LTcaptype{none} % do not increment counter
\begin{longtable}[]{@{}
  >{\raggedright\arraybackslash}p{(\linewidth - 12\tabcolsep) * \real{0.0661}}
  >{\raggedright\arraybackslash}p{(\linewidth - 12\tabcolsep) * \real{0.1074}}
  >{\raggedright\arraybackslash}p{(\linewidth - 12\tabcolsep) * \real{0.1736}}
  >{\raggedright\arraybackslash}p{(\linewidth - 12\tabcolsep) * \real{0.1570}}
  >{\raggedright\arraybackslash}p{(\linewidth - 12\tabcolsep) * \real{0.0909}}
  >{\raggedright\arraybackslash}p{(\linewidth - 12\tabcolsep) * \real{0.2479}}
  >{\raggedright\arraybackslash}p{(\linewidth - 12\tabcolsep) * \real{0.1570}}@{}}
\toprule\noalign{}
\begin{minipage}[b]{\linewidth}\raggedright
\textbf{\emph{ID}}
\end{minipage} & \begin{minipage}[b]{\linewidth}\raggedright
\textbf{Categoría}
\end{minipage} & \begin{minipage}[b]{\linewidth}\raggedright
\textbf{Vulnerabilidad}
\end{minipage} & \begin{minipage}[b]{\linewidth}\raggedright
\textbf{Tipo}
\end{minipage} & \begin{minipage}[b]{\linewidth}\raggedright
\textbf{Impacto}
\end{minipage} & \begin{minipage}[b]{\linewidth}\raggedright
\textbf{Detectable automáticamente}
\end{minipage} & \begin{minipage}[b]{\linewidth}\raggedright
\textbf{Ejemplo sección}
\end{minipage} \\
\midrule\noalign{}
\endhead
\bottomrule\noalign{}
\endlastfoot
\emph{V1} & Técnica & Reentrancy & Ejecución & Crítico & Sí
(Slither/Mythril) & ANEXO A.1.1 \\
\emph{V2} & Técnica & Overflow & Aritmético & Medio & Sí (Slither) &
ANEXO A.1.2 \\
\emph{V3} & Técnica & Delegatecall & Ejecución & Crítico & Parcial &
ANEXO A.1.3 \\
\emph{V4} & Control & Acceso no restringido & Autorización & Crítico &
Sí & ANEXO A.2.1 \\
\emph{V5} & Control & tx.origin & Autenticación & Alto & Sí & ANEXO
A.2.2 \\
\emph{V6} & Económica & Front-running & MEV & Alto & No & ANEXO A.3.1 \\
\emph{V7} & Económica & Oráculos & Dependencia externa & Crítico & No &
ANEXO A.3.2 \\
\emph{V8} & Lógica & Error de balance & Lógica & Variable & No & ANEXO
A.4.1 \\
\end{longtable}
}
